\section{Analysis of Empirical, Finite-Width Dynamics}\label{section:finite_sample_analysis}

The first step is to set up a coupling between the finite-width empirical dynamics and the infinite-width population dynamics --- we will eventually show that these two dynamics do not diverge very far, even up to time $\Tstar$ (defined in \Cref{thm:pop_main_thm}).

\paragraph{Coupling Between Empirical and Population Dynamics.} To analyze the difference between the empirical dynamics (resulting in $\rhohat_t$) and the population dynamics (resulting in $\rho_t$) we define an intermediate process $\rhobar_t$, which has finite support. Specifically, $\rhobar_0 = \rhohat_0$, and $\rhobar_t$ will always have support size $m$, but the particles in $\rhobar$ are then updated according to the infinite-width population dynamics. We formally define $\rhobar_t$ as follows:
\begin{align}
\rhobar_{0} & = \rhohat_0 = \textup{unif}(\{\chi_1,\dots, \chi_m\}) \comma \nonumber\\
\text{ and } \rhobar_t & = \textup{distribution of $\u_t(\chi)$ (where $\chi \sim \rhobar_0$)}  \period \label{eqn:13}
\end{align}
In other words, $\rhobar_t$ consists of the particles in $\rho_t$ whose initialization is in $\{\chi_1, \ldots, \chi_m\}$.

Now, we define a coupling between $\rhobar_t$ and $\rhohat_t$. For $\chi \sim \textup{unif}(\{\chi_1,\dots, \chi_m\})$, let $\ubar_t(\chi) = \u_t(\chi)$. Let $\coup_t$ be the joint distribution of $(\ubar_t(\chi), \uhat_t(\chi))$. Then, $\coup_t$ forms a natural coupling between $\rhobar_t$ and $\rhohat_t$. We will use $\ubar_t$ and $\uhat_t$ as shorthands for the random variables  $\ubar_t(\chi), \uhat_t(\chi)$ respectively in the rest of this section. Additionally, we define the average distance $\deltaBar_t^2 := \E_{(\uhat_t, \ubar_t)\sim \coup_t} [\|\uhat_t - \ubar_t\|_2^2]$. Intuitively, $f_{\rho_t}(x)$ and $f_{\rhohat_t}(x)$ are close when $\uhat_t$ and $\ubar_t$ are close, which is formalized by the following lemma:

\begin{lemma}\label{lemma:finite_sample_width_main_lemma}
In the setting of \Cref{thm:empirical-main}, suppose the network width $m$ satisfies $m \leq d^{O(\log d)}$, and let $T \leq \frac{d^C}{\sigmaCoeffTwo^2 + \sigmaCoeffFour^2}$ for some universal constant $C > 0$. Then, with probability at least $1 - \exp(-d^2)$ over the randomness of $\{\chi_1, \ldots, \chi_m\}$, we have for all $t \leq T$ that
\begin{align}
\E_{x \sim \bbS^{d - 1}} [(f_{\rho_t}(x) - f_{\rhobar_t}(x))^2]
& \lesssim \frac{(\sigmaCoeffTwo^2 + \sigmaCoeffFour^2) d^2 (\log d)^{O(1)}}{m} \period
\end{align}
Additionally, for all $t \geq 0$, we have
\begin{align}
\E_{x \sim \bbS^{d - 1}} [(f_{\rhohat_t}(x) - f_{\rhobar_t}(x))^2]
& \lesssim (\sigmaCoeffTwo^2 + \sigmaCoeffFour^2) \deltaBar_t^2 \period
\end{align}
\end{lemma}

The proof of \Cref{lemma:finite_sample_width_main_lemma} is deferred to \Cref{section:proof_finite_sample}. As a simple corollary of \Cref{lemma:finite_sample_width_main_lemma}, we can upper bound $\Exp_{x\sim\Sp}[(f_{\rho_t}(x)-f_{\rhohat_t}(x))^2]$ by the triangle inequality. Thus, so long as $\deltaBar_{\Tstar}$ is small, we can show that the empirical and population dynamics achieve similar test error at time $\Tstar$.

\paragraph{Upper Bound for $\deltaBar_{\Tstar}$.} The following lemma gives our upper bound on $\deltaBar_{\Tstar}$:

\begin{lemma}[Final Bound on $\deltaBar_{\Tstar}$] 
	\label{lemma:final_bound_Delta_avg}
     In the setting of Theorem~\ref{thm:empirical-main}, we have
	\begin{align}
		\deltaBar_{\Tstar}\le d^{-\frac{\samplebound-1}{4}} \period
	\end{align}	
\end{lemma}



We prove \Cref{lemma:final_bound_Delta_avg} by induction over the time $t \le \Tstar$.
Let us define the maximal distance
$\deltaMaxt{t} := \max_{(\uhat_t, \ubar_t) \in \supp(\coup_t)}\|\uhat_t - \ubar_t\|_2$ where the max is over the \textit{finite} support of the coupling $\coup_t$. 
We will control $\deltaBar_t$ and $\deltaMaxt{t}$ simultaneously using induction. The inductive hypothesis is: 
\begin{assumption}[Inductive Hypothesis] \label{assumption:inductive_hypothesis}
	For some universal constant $1 > \avgbound > 1/2$ and $\maxbound \in (0, \avgbound-\frac{1}{2})$, we say that the inductive hypothesis holds at time $T$ if, for all $0\le t\leq T$, we have 
	\begin{align}
	\deltaBar_t \leq d^{-\avgbound}~~ \textup{ and } ~~\deltaMaxt{t} \leq d^{-\maxbound} \period
	\end{align}
\end{assumption}
Showing that the inductive hypothesis holds for all $t$ requires studying the growth of the error $\norm{\uhat_t - \ubar_t}$. We analyze it using the following decomposition:~\footnotemark \footnotetext{Note that $\uhat_t-\ubar_t$ and $A_t, B_t$ and $C_t$ implicitly depend on the initialization $\chi$, and could be written as $A_t(\chi), B_t(\chi), C_t(\chi)$, and $\delta_t(\chi)$, but we omit this dependence for ease of presentation.}
\begin{align}
	\frac{d}{dt} \|\uhat_t-\ubar_t\|_2^2 =& A_t + B_t + C_t \comma
\end{align}
where $A_t := - 2 \langle \pgrad_{\uhat} L(\rho_t) - \pgrad_{\ubar} L(\rho_t), \uhat_t - \ubar_t \rangle$,
$B_t := -2 \langle \pgrad_{\uhat} L(\rhohat_t) - \pgrad_{\uhat} L(\rho_t), \uhat_t - \ubar_t \rangle$, and $C_t = -2 \langle \pgrad_{\uhat} \Lhat(\rhohat_t) - \pgrad_{\uhat} L(\rhohat_t), \uhat_t - \ubar_t \rangle.$ Intuitively, $A_t$ captures the growth of the coupling error due to the population dynamics, $B_t$ captures the growth of the coupling error due to the discrepancy between the gradients from the finite-width and infinite-width networks, and $C_t$ captures the growth of the coupling error due to the discrepancy between finite samples and infinite samples. 

We will use the following lemma to give an upper bound on each of these three terms.

\begin{lemma}[Bounds on $A_t, B_t, C_t$] \label{lemma:final_bounds_ABC}
In the setting of \Cref{thm:empirical-main}, suppose the inductive hypothesis (\Cref{assumption:inductive_hypothesis}) holds up to time $t$, for some $t \leq \Tstar$. Let $\vv_t := \uhat_t - \ubar_t$. Let $T_1$ be the runtime of Phase 1 (where Phase 1 is defined in \Cref{def:phase_1}). Then, we have
\begin{align}
    A_t &\le \begin{cases} 4\sigmaCoeffTwo^2|D_{2,t}|\norm{\vv_t}^2 + O\big( \frac{\sigmaCoeffTwo^2}{\log d}|D_{2,t}| \norm{\vv_t}^2 \big) & \text{if } 0\le t \le T_1 \\
    O(1)\cdot (\sigmaCoeffTwo^2 + \sigmaCoeffFour^2)\gamma_2 \norm{\vv_t}^2 & \text{if } T_1\le t \le \Tstar 
\end{cases} \period\label{eqn:at_bound}
\end{align}
Additionally, assuming that the width $m$ is at most $d^{O(\log d)}$, we have with probability at least $1 - e^{-d^2}$ over the randomness of the initialization $\rhohat_0$ that $B_t \lesssim (\sigmaCoeffTwo^2 + \sigmaCoeffFour^2) \cdot \Big(\frac{d (\log d)^{O(1)}}{\sqrt{m}} + \deltaBar_t \Big) \|\delta_t\|_2$ and $\E[B_t] \lesssim \frac{(\sigmaCoeffTwo^2 + \sigmaCoeffFour^2) d (\log d)^{O(1)}}{\sqrt{m}} \deltaBar_t + (\sigmaCoeffTwo^2 + \sigmaCoeffFour^2) \deltaBar_t^3$. 

Finally, assume that the number of samples $n \leq d^C$ and the width $m$ is at most $d^C$ for any universal constant $C > 0$. Assume that $\deltaBar_t \leq \frac{1}{\sqrt{d}}$ (i.e. $\avgbound > \frac{1}{2}$) and that $\maxbound < \frac{1}{2}$. Then, with probability $1 - \frac{1}{d^{\Omega(\log d)}}$ over the randomness of the dataset, we have $C_t \lesssim (\sigmaCoeffTwo^2 + \sigmaCoeffFour^2) (\log d)^{O(1)} \cdot \Big(\sqrt{\frac{d}{n}} \|\delta_t\|_2 + \frac{d^{2 - 2\maxbound}}{n} \deltaBar \|\delta_t\|_2^2 + \frac{d^{2.5 - 2\maxbound}}{n} \deltaBar^2 \|\delta_t\|_2  + \frac{d^{4 - 4\maxbound}}{n} \deltaBar^2 \|\delta_t\|_2^2 \Big)$.
\end{lemma}

The proof can be found in \Cref{section:proof_finite_sample}. We give an explanation for each of these bounds below.

For $A_t$, we establish two separate bounds, one which holds during Phase 1, and one which holds during Phases 2 and 3. Intuitively, the signal part in a neuron (i.e. $\w^2$) grows at a rate $4 \sigmaCoeffTwo^2 |D_{2, t}|$ during Phase 1 (which follows from \Cref{lem:1-d-traj}) and in \Cref{eqn:at_bound} we show that the growth rate of the coupling error due to the growth in the signal is also at most $4 \sigmaCoeffTwo^2 |D_{2, t}|$. Intuitively, the growth rates match because the signal parts of the neurons follow similar dynamics to a power method during Phase 1. For example, in the worst case when $\delta_t$ is mostly in the direction of $\e$, the factor by which $\delta_t$ grows is the same as that of $\w$. More mathematically, this is due to the fact that the first-order term in \Cref{lem:1-d-traj} is the dominant term in $v(\w) - v(\what)$ for any two particles $\w, \what$ (where $v(\w)$ is the one-dimensional velocity defined in the previous section). To get a sample complexity better than NTK, the precise constant factor $4$ in the growth rate $4 \sigmaCoeffTwo^2 |D_{2, t}| \w$ is important --- a larger constant would lead to $\delta_{T_1}$ being larger by a $\poly(d)$ factor, thus increasing the sample complexity by $\poly(d)$.

The same bound for $A_t$ no longer applies during Phases 2 and 3. This is because the dynamics of $\w$ are no longer similar to a power method, and the higher-order terms may make a larger contribution to $v(\w) - v(\what)$ than to the growth of $\w$, or vice versa. Thus we use a looser bound $O(1)\cdot (\sigmaCoeffTwo^2 + \sigmaCoeffFour^2)\gamma_2 \norm{\vv_t}^2$ in Eq.~\eqref{eqn:at_bound}. However, since the remaining running time after Phase 1, $\Tstar-T_1$, is very short (only $\poly(\log\log d)$ --- this corresponds to the second term of the running time in \Cref{thm:pop_main_thm}), the total growth of the coupling error is at most $\exp(\poly(\log\log d))$, which is sub-polynomial and does not contribute to the sample complexity. Note that if the running time during Phases 2 and 3 is longer (e.g. $O(\log d)$) then the coupling error would grow by an additional $d^{O(1)}$ factor during Phases 2 and 3, which would make our sample complexity worse than NTK.

Our upper bounds on $B_t$ and $C_t$ capture the growth of coupling error due to the finite width and samples. We establish a stronger bound for $\E[B_t]$ than for $B_t$. In our proof, we use the bounds for $B_t$ and $\Exp[B_t]$ to prove the inductive hypothesis for $\deltaBar_t$ and $\deltaMaxt{t}$ respectively. We prove the bound for $C_t$ by expanding the error due to finite samples into second-order and fourth-order polynomials and applying concentration inequalities for higher moments (\Cref{lemma:main_concentration_lemma}).

All of these bounds together control the growth of $\norm{\uhat_t - \ubar_t}$ at time $t$. Using Lemma~\ref{lemma:final_bounds_ABC} we can show that for some properly chosen $\avgbound$ and $\maxbound$, the inductive hypothesis in \Cref{assumption:inductive_hypothesis} holds until $\Tstar$ (see Lemma~\ref{lemma:IH_is_maintained} for the statement), which naturally leads to the bound in Lemma~\ref{lemma:final_bound_Delta_avg}. The full proof can be found in Appendix~\ref{section:proof_finite_sample}. 